\documentclass[11pt,a4paper]{article}
\usepackage[margin=2.5cm]{geometry}
\usepackage{amsmath,amssymb}
\usepackage{booktabs}
\usepackage{longtable}
\usepackage{array}
\usepackage[hidelinks]{hyperref}
\usepackage{xcolor}
\emergencystretch=2em

\newcommand{\vu}{\mathbf{u}}
\newcommand{\dd}{\,\mathrm{d}}
\newcommand{\code}[1]{\texttt{#1}}
\newcommand{\Yf}{Y_{\!F}}
\newcommand{\YO}{Y_{O_2}}

\title{Coupled atmosphere--fire model: equations, boundary conditions\\and numerical methods of the PyTorch implementation}
\author{Documentation of \code{src/} (branch \code{more-physics})}
\date{November 2026}

\begin{document}
\maketitle

\begin{abstract}
This document describes the model solved by the PyTorch code in \code{src/}. Two variants share
one flow solver: the \emph{published model} of San Mart\'in \& Torres (2025), with a single
temperature for gas and fuel and radiation as diffusion, and an \emph{extended model} with a
separate solid fuel temperature, fuel moisture, heat-driven pyrolysis, transported fuel gas and
oxygen with mixing-limited combustion, and radiation transport. Section~\ref{sec:numerics}
describes the discretization and solvers, Section~\ref{sec:params} lists every input parameter.
Status: the extended model is verified for consistency (analytic radiation test, energy budget,
equality of 2D runs and $y$-invariant 3D runs, exact reproduction of the original code in legacy
mode) but not yet calibrated or validated against the CSIRO F19 experiment.
\end{abstract}

\tableofcontents

%======================================================================
\section{Overview}
%======================================================================

\paragraph{Model variants.}
\begin{itemize}
  \item \textbf{Published model}: \code{[fuel] solid\_temperature = False},
        \code{[radiation] model = diffusion}. The option \code{--legacy} additionally restores the
        original numerics (central temperature advection, explicit reaction, original sink term,
        original pressure scaling), reproducing \code{src\_legacy/} to round-off.
  \item \textbf{Extended model} (default): sections~\ref{sec:solid}, \ref{sec:combustion} and
        \ref{sec:radiation}.
\end{itemize}

\paragraph{Notation.} Velocity $\vu = (u,[v,]w)$ with $w$ vertical; in 2D output files the
vertical velocity is stored as \code{v} and the vertical coordinate as \code{y}. Gas temperature
$T$; fuel fraction $Y\in[0,1]$, the fraction of the initial dry fuel left; gravity
$\mathbf{g}=(0,[0,]g)$ with $g=-9.807$~m\,s$^{-2}$. Arrays are stored as $(z,x)$ in 2D and
$(z,y,x)$ in 3D.

\paragraph{State variables.}
\begin{center}
\begin{tabular}{lll}
\toprule
Symbol & Field (output name) & Model \\
\midrule
$\vu$ & velocity (\code{u}, \code{v}[, \code{w}]) & both \\
$T$ & gas temperature (\code{T}) & both \\
$Y$ & fuel fraction left in the bed (\code{Y}) & both \\
$T_s$ & solid fuel temperature (\code{Ts}) & extended \\
$W$ & fuel moisture, kg water per kg dry fuel (\code{W}) & extended \\
$\Yf$ & mass fraction of pyrolysis gas (\code{YF}) & extended \\
$\YO$ & mass fraction of oxygen (\code{YO2}) & extended \\
$p$ & hydrodynamic pressure (\code{p}) & both \\
\bottomrule
\end{tabular}
\end{center}

%======================================================================
\section{Gas phase}
%======================================================================

\subsection{Low Mach number approximation}
The thermodynamic pressure is constant, so the ideal gas law gives
\begin{equation}
  \rho = \frac{\rho_\infty T_\infty}{T}
\end{equation}
(with \code{[density] constant = True}, $\rho=\rho_\infty$). Mass conservation becomes a constraint
on the velocity divergence:
\begin{equation}\label{eq:div}
  \nabla\cdot\vu = D \equiv \frac{\dot q}{\rho c_p T} + \frac{\dot m'''}{\rho},
\end{equation}
where $\dot q$ (W\,m$^{-3}$) is the net heating of the gas (all terms of the energy equation except
advection) and $\dot m'''$ (kg\,m$^{-3}$\,s$^{-1}$) the gas produced by pyrolysis (extended model).

\subsection{Momentum}
\begin{equation}\label{eq:momentum}
  \frac{\partial\vu}{\partial t} + (\vu\cdot\nabla)\vu
  = -\frac{1}{\rho}\nabla p + \nu\nabla^2\vu + \frac{\nu}{3}\nabla D
  + \frac{\rho-\rho_\infty}{\rho}\,\mathbf{g} - \mathbf{F}_d - \boldsymbol{\tau}_{sgs},
\end{equation}
with $\nu=\mu/\rho$, the imposed divergence $D$ of \eqref{eq:div}, the drag of the fuel elements
\begin{equation}
  \mathbf{F}_d = \tfrac12 C_d\,\alpha_s\sigma_s\,Y\,|\vu|\,\vu,
\end{equation}
and the subgrid stresses $\boldsymbol{\tau}_{sgs}$ of Section~\ref{sec:turbulence}.

\subsection{Energy}
The specific heat $c_p$ is constant by default; with \code{[gas] cp\_model = variable} it is the
temperature-dependent value of air from ideal-gas tables (1005~J\,kg$^{-1}$K$^{-1}$ at 300~K, 1142 at
1000~K, 1253 at 2000~K, interpolated linearly), used in the energy equation, the low-Mach constraint
\eqref{eq:div} and the energy budget. With \code{cp\_model = mixture} (extended model) it is weighted
by composition, $c_p=\sum_kY_kc_{p,k}(T)$ over N$_2$, O$_2$, CO$_2$ and H$_2$O (rounded JANAF values). The
products are inferred from the oxygen consumed: the deficit $Y_{O_2,\infty}(1-\Yf)-\YO$ corresponds to
$(1+1/s)$~kg of products per kg of oxygen, split into CO$_2$ and H$_2$O by the cellulose stoichiometry
C$_6$H$_{10}$O$_5$ + 6\,O$_2$ $\to$ 6\,CO$_2$ + 5\,H$_2$O; the fuel gas and the rest are counted as N$_2$.
Fully burned gas has $c_p\approx1480$~J\,kg$^{-1}$K$^{-1}$ at 2000~K, 18\% more than air.

Transport properties are constant by default. With \code{[gas] transport = sutherland} the viscosity
and the conductivity follow Sutherland's law, $\mu(T)=\mu_\infty(T/T_\infty)^{3/2}(T_\infty+S)/(T+S)$,
with $S=110.4$~K for $\mu$ and 194~K for $\kappa$, scaled to the constant values at $T_\infty$. The
conduction term is then $\nabla\cdot(\kappa\nabla T)=\kappa\nabla^2T+\kappa'|\nabla T|^2$, the gradient of
$\mu$ is neglected in the momentum equation, and the convective heat transfer to the fuel uses the
gas properties at the film temperature $(T+T_s)/2$. Molecular transport is several orders of
magnitude smaller than the subgrid transport at the resolutions used, so its main effect is on the
heat transfer to the fuel.
\begin{equation}\label{eq:energy}
  \frac{\partial T}{\partial t} + \vu\cdot\nabla T = \frac{\dot q}{\rho c_p} - q_{sgs}(T).
\end{equation}

\paragraph{Published model.}
\begin{equation}
  \dot q = \nabla\cdot\big(k(T)\nabla T\big) + \rho c_p\,q(T,Y),\qquad k(T)=\kappa+4\sigma\delta T^3,
\end{equation}
where $4\sigma\delta T^3$ represents radiation as diffusion (optical path length $\delta$) and
\begin{equation}
  q = \frac{H_C\,Y\,R(T)}{c_p} + q_{sink}
\end{equation}
with the reaction rate $R$ of \eqref{eq:arrhenius} and the convective exchange $q_{sink}$ of
Section~\ref{sec:sink}.

\paragraph{Extended model.}
\begin{equation}\label{eq:qdot-ext}
  \dot q = \kappa\nabla^2T + \kappa_g\,(G-4\sigma T^4) + \dot\omega_F H_C - Q_{conv},
\end{equation}
with molecular conduction, radiation absorbed minus emitted by the gas (incident radiation $G$,
Section~\ref{sec:radiation}; gray absorption coefficient $\kappa_g$), the heat of gas-phase
combustion (Section~\ref{sec:combustion}) and the convective heat transferred to the fuel
(Section~\ref{sec:solid}). With \code{[radiation] model = diffusion} the first two terms are
replaced by $\nabla\cdot(k(T)\nabla T)$.

%======================================================================
\section{Solid fuel}\label{sec:solid}
%======================================================================

\subsection{Fuel bed}
The fuel occupies the nodes up to a height $Y_h$ above the terrain (optionally limited
horizontally), with $Y=1$ initially. The elements are characterized by the solid volume fraction
$\alpha_s$, the surface-to-volume ratio $\sigma_s$ (equivalent cylinder diameter $d=4/\sigma_s$) and,
in the extended model, the dry bulk density $\rho_b$. Fuel surface per unit volume:
$a=\alpha_s\sigma_s Y$.

\subsection{Fuel temperature and moisture (extended model)}
\begin{equation}\label{eq:solid}
  \rho_b\,(Y c_s + W c_w)\,\frac{\dd T_s}{\dd t} = Q_{conv} + Q_{rad} - \dot m''' H_{pyr},
\end{equation}
\begin{align}
  Q_{conv} &= h\,a\,(T-T_s), \qquad
  h = \max\!\Big(h_c|T-T_s|^{1/3},\ \frac{\kappa}{d}\,0.683\,Re^{0.466}Pr^{1/3}\Big),
  \quad Re=\frac{|\vu|\,d\,\rho}{\mu},\\
  Q_{rad} &= \frac{a}{4}\,\big(G-4\sigma T_s^4\big),
\end{align}
i.e.\ the larger of a natural-convection estimate and the forced convection of a cylinder
(Hilpert correlation), and the radiation absorbed by randomly oriented elements (absorption
coefficient $a/4$). Wet fuel ($W>0$) at $T_{boil}$ keeps its temperature and evaporates:
\begin{equation}
  \rho_b\frac{\dd W}{\dd t} = -\frac{\max(Q_{conv}+Q_{rad},0)}{L_v}.
\end{equation}
Water vapor is not tracked. Only dry fuel pyrolyses. Where the fuel is used up
($Y\le10^{-3}$, $W=0$), $T_s$ is set equal to $T$.

\subsection{Pyrolysis (extended model)}
\paragraph{Linear (default, after Mell et al.\ 2007).} Dry fuel converts the fraction
\begin{equation}
  f = \min\!\Big(\max\!\Big(\frac{T_s-T_{pyr}}{T_{pyr,end}-T_{pyr}},0\Big),1\Big)
\end{equation}
of the heat it receives into pyrolysis:
\begin{equation}
  \dot m''' = \frac{f\,\max(Q_{conv}+Q_{rad},0)}{H_{pyr}},\qquad \rho_b\frac{\dd Y}{\dd t}=-\dot m'''.
\end{equation}
Above $T_{pyr,end}$ all net heat goes to pyrolysis, so $T_s$ stays near $T_{pyr,end}$ while fuel is left.

\paragraph{Arrhenius.} $\dot m''' = \rho_b\,Y_f\,R(T_s)\,Y$ with $R$ of \eqref{eq:arrhenius} at the
fuel temperature.

In both cases $\dot m'''\le\rho_bY/\Delta t$ (no more fuel than is present).

\subsection{Reaction rate of the published model}
\begin{equation}\label{eq:arrhenius}
  \frac{\dd Y}{\dd t} = -Y_f\,R(T)\,Y,\qquad R(T)=A\,e^{-T_{act}/T}\,H(T-T_{pc}),
\end{equation}
with the Heaviside step $H$. $Y_f$ sets how fast the fuel is consumed and the total heat per unit
of fuel, $H_C/(c_pY_f)$ in temperature units, but not the instantaneous heat release.

\subsection{Convective exchange of the published model}\label{sec:sink}
With \code{sink = surface} (default) the top fuel node of each column (highest node with
$Y>10^{-3}$) exchanges the flux per unit ground area
\begin{equation}
  Q = h\,(T_g-T_f),\qquad h=h_c\Big(\frac{|T_g-T_f|}{Y_h}\Big)^{1/4},
\end{equation}
with the gas node above it; $Q/(\rho c_p\Delta z)$ is added to the fuel node and subtracted from
the gas node (energy conserving, grid independent). \code{sink = legacy} keeps the original term.

%======================================================================
\section{Gas-phase combustion (extended model)}\label{sec:combustion}
%======================================================================
\begin{align}
  \frac{\partial\Yf}{\partial t} + \vu\cdot\nabla\Yf &= \mathcal{D}\nabla^2\Yf - q_{sgs}(\Yf)
    + \frac{\dot m'''(1-\Yf) - \dot\omega_F}{\rho},\\
  \frac{\partial\YO}{\partial t} + \vu\cdot\nabla\YO &= \mathcal{D}\nabla^2\YO - q_{sgs}(\YO)
    - \frac{s\,\dot\omega_F + \dot m'''\YO}{\rho},
\end{align}
with $\mathcal{D}=\kappa/(\rho c_p)$ (unit Lewis number) and turbulent Schmidt number equal to $Pr$.
The terms with $\dot m'''$ account for the added mass of pyrolysis gas. Combustion is mixing limited:
\begin{equation}
  \dot\omega_F = \frac{\rho\,\min(\Yf,\ \YO/s)}{\tau_{mix}},\qquad \tau_{mix}=\sqrt{2\Delta/|g|}
\end{equation}
(the FDS buoyancy time scale, $\Delta$ the filter width) unless \code{tau\_mix} is given; $s$ is the
stoichiometric oxygen-to-fuel mass ratio. Products are not tracked individually.

%======================================================================
\section{Turbulence}\label{sec:turbulence}
%======================================================================
Smagorinsky model with Van Driest damping:
\begin{equation}
  \mu_t = \rho\,(C_s\Delta f_w)^2|S|,\qquad f_w = 1-e^{-z^+/26},
\end{equation}
\begin{equation}
  |S|^2 = 2\sum_i(\partial_iu_i)^2 + \sum_{i<j}(\partial_ju_i+\partial_iu_j)^2 - \tfrac23(\nabla\cdot\vu)^2,
\end{equation}
$z^+=d_w\rho u_\tau/\mu$ with $d_w$ the height above the terrain, $u_\tau=\sqrt{\tau_w/\rho}$ and
$\tau_w=\mu|\partial_zu+\partial_xw|$ at the ground (plus the $y$ component in 3D). Filter width
$\Delta=(\Delta x\Delta z)^{1/2}$ (2D), $(\Delta x\Delta y\Delta z)^{1/3}$ (3D), or a fixed width
(\code{filter\_width}). With the grid-based width the mixing length shrinks under refinement and
the results do not converge; tests on the F19 case showed a fixed width removes this.
\begin{align}
  \tau_{sgs,i} &= -\frac{2}{\rho}\big(\nabla\mu_t\cdot S_{i\cdot} + \mu_t\,\nabla\cdot S_{i\cdot}\big),\\
  q_{sgs}(\phi) &= -\frac{1}{Pr\,\rho}\big(\nabla\mu_t\cdot\nabla\phi + \mu_t\nabla^2\phi\big),
\end{align}
with $S$ the deviatoric strain rate and $\phi\in\{T,\Yf,\YO\}$.

%======================================================================
\section{Radiation (extended model)}\label{sec:radiation}
%======================================================================
Gray radiative transfer:
\begin{equation}
  \boldsymbol{\Omega}\cdot\nabla I = \kappa\,(I_b-I),\qquad \kappa=\kappa_g+\kappa_s,\quad
  \kappa_s=\frac{\alpha_s\sigma_sY}{4},\qquad
  I_b = \frac{\sigma(\kappa_gT^4+\kappa_sT_s^4)}{\pi\kappa},
\end{equation}
and $G=\int_{4\pi}I\dd\Omega$. In the single-temperature model with radiation transport the fuel
emits at $T$. Boundary conditions: black ground at the temperature of the lowest gas row; ambient
radiation $\sigma T_\infty^4/\pi$ entering at the top; periodic horizontally; terrain nodes opaque
at their own temperature. At uniform $T_\infty$, $G=4\sigma T_\infty^4$ and there is no net exchange.

%======================================================================
\section{Initial and boundary conditions}\label{sec:ibc}
%======================================================================

\subsection{Initial conditions}
\begin{itemize}
  \item Wind (\code{[wind] shape}), with $d_w$ the height above the terrain; other components zero:
    \begin{itemize}
      \item \code{power law} (default): $u=u_r(d_w/z_r)^{1/7}$;
      \item \code{monin-obukhov}: Monin--Obukhov similarity theory, as in the FDS \code{\&WIND} input,
            \begin{equation}
              u = \frac{u_*}{\kappa}\Big[\ln\frac{d_w}{z_0} - \psi_m\Big(\frac{d_w}{L}\Big) + \psi_m\Big(\frac{z_0}{L}\Big)\Big],
            \end{equation}
            with the roughness length $z_0$ (\code{z\_0}), the Obukhov length $L$ (\code{L}; omitted: neutral),
            $\kappa=0.41$, zero below $z_0$, and $u_*$ chosen so that $u(z_r)=u_r$. Stability function
            (Businger--Dyer, integrated by Paulson 1970): for $\zeta=d_w/L<0$, with $x=(1-16\zeta)^{1/4}$,
            $\psi_m=2\ln\frac{1+x}{2}+\ln\frac{1+x^2}{2}-2\arctan x+\frac{\pi}{2}$; for $\zeta\ge0$,
            $\psi_m=-5\zeta$. The FDS Case F19 values are $u_r=4.9$~m\,s$^{-1}$ at $z_r=2$~m, $z_0=0.03$~m,
            $L=-500$~m ($u_*=0.48$~m\,s$^{-1}$). Near the ground this profile is much weaker than the
            power law: 1.4, 2.5 and 3.3~m\,s$^{-1}$ at 0.1, 0.25 and 0.51~m, against 3.1, 3.6 and
            4.0~m\,s$^{-1}$. Only the velocity is initialized from the similarity theory; the temperature
            stays uniform.
      \item \code{log} (neutral logarithmic profile from $u_*$) or \code{constant}.
    \end{itemize}
  \item Fuel: $Y=1$ in the bed, $0$ elsewhere.
  \item Gas temperature: $T=T_\infty+(T_{hot}-T_\infty)\,S$ with the ignition shape $S$:
    \begin{itemize}
      \item \code{gaussian} (default): $S=\exp\!\big(-(x-x_c)^2/L^2 - z^2/H^2\big)$,
            $x_c=(x_{start}+x_{end})/2$, $L=x_{end}-x_{start}$, $H=z_{end}-z_{start}$ (3D: also Gaussian in $y$);
      \item \code{plate}: indicator of the box $[x_{start},x_{end}]\times[z_{start},z_{end}]$ (3D: and $[y_{start},y_{end}]$);
      \item \code{line} (3D): the 2D Gaussian, uniform for $y_{start}\le y\le y_{end}$.
    \end{itemize}
  \item Extended model: $T_s=\min(T,T_{pyr,end})$ (linear pyrolysis) or $T_s=T$; $W=$\,\code{moisture}$\cdot Y$
        except where $T\ge T_{boil}$ (dry); $\Yf=0$; $\YO=0.232$.
  \item Optional sustained source: for $t\le$\,\code{t\_source}, $T$ is held at its initial value
        where $S>0.01$.
\end{itemize}

\subsection{Boundary conditions}
\begin{center}
\begin{tabular}{p{2.6cm}p{4.2cm}p{2.6cm}p{2.3cm}p{1.8cm}}
\toprule
Boundary & Velocity & $T$ & $\Yf$, $\YO$ & $p$ \\
\midrule
Horizontal & periodic & periodic & periodic & periodic \\
Ground $z_{\min}$ & Dirichlet: initial profile at the ground (zero for the power law, i.e.\ no slip) & zero gradient (adiabatic) & zero gradient & zero gradient \\
Top $z_{\max}$ & zero gradient & zero gradient & zero gradient & $p=0$ \\
\bottomrule
\end{tabular}
\end{center}
Discretization of the zero-gradient conditions: velocities use the second-order extrapolation
$\phi_0=(4\phi_1-\phi_2)/3$; the scalars ($T$, $\Yf$, $\YO$) use the first-order form $\phi_0=\phi_1$.
The second-order form leaves the range of the interior values where the profile is curved: at the
fire front, with a hot layer just above cold ground ($T_1=348$~K, $T_2=584$~K), it gives 270~K, below
ambient, and the temperature bound then added energy at every step (up to about 18\% of the heat of
combustion over the first 6~s of the F19 case at $\Delta x=1.17$~m). The published code, reproduced
with \code{--legacy}, uses the second-order form for all variables.
$Y$, $T_s$ and $W$ are not transported and need no boundary condition.

\subsection{Wind forcing}
In a periodic domain nothing drives the wind: ground friction drains its momentum and the
near-ground wind decays. Without fire, on a 4.7~m~$\times$~0.31~m grid, the horizontal-mean wind at
2~m fell from 4.8 to 2.2~m\,s$^{-1}$ in 60~s, and at 0.3~m from 2.7 to 0.6~m\,s$^{-1}$; only the upper
layers kept their speed. With \code{[wind] forcing = profile} a body force relaxes the horizontal-mean
velocity at each height toward the initial profile,
\begin{equation}
  \mathbf{f}(z) = \frac{\overline{\vu}_0(z) - \overline{\vu}(z,t)}{\tau_f},
\end{equation}
where the overbar is the average over the fluid nodes of a horizontal plane and $\tau_f$ is
\code{forcing\_time} (default 1~s). The force is uniform in the horizontal, so it does not distort
the local flow around the fire; with a large domain the fire changes the plane averages only
slightly. In the test above, $\tau_f=0.5$~s held the wind at 2~m and above within 1\%, but with a
no-slip ground the lowest cells still lost 18--25\%: on coarse grids the no-slip condition
overestimates the wall friction. The wall model below removes this deficit.

\subsection{Wall model}
With \code{[wall] model = log} the no-slip condition at the ground (and on terrain surfaces) is
replaced by a wall-stress model. The friction velocity follows from the log law at the first node
above the ground, at height $z_1$, with the horizontal velocity $\mathbf{u}_1$:
\begin{equation}
  u_\tau = \frac{\kappa\,|\mathbf{u}_1|}{\ln(z_1/z_0) - \psi_m(z_1/L) + \psi_m(z_0/L)},
\end{equation}
with the stability function of the wind profile when an Obukhov length is given, and the roughness
length $z_0$ (\code{[wall] z\_0}, default the wind's \code{z\_0}). The tangential velocity at the wall
is set equal to that of the first node, so the resolved stress at the wall vanishes, the vertical
velocity at the wall is zero, and the wall stress enters as a momentum sink in the first cell,
$-u_\tau^2\,u_i/|\mathbf{u}_1|/\Delta z$. The Van Driest damping of the turbulence model uses the modeled
$u_\tau$. Building walls keep the no-slip condition.

Verification (no fire, Monin--Obukhov profile of the F19 case, 4.7~m~$\times$~0.31~m grid, 60~s):
with the wall model and forcing ($\tau_f=1$~s) the horizontal-mean wind at 0.3, 0.6, 2 and 10~m
was 2.68, 3.51, 4.79 and 6.72~m\,s$^{-1}$, against the initial 2.74, 3.55, 4.83 and 6.72 (within 2\%),
and the modeled friction velocity was about 0.47~m\,s$^{-1}$, against 0.48 for the profile. With
no-slip and the same forcing the lowest values were 1.85 and 2.60; with neither, 0.56 and 0.76.

\subsection{Buildings}
Buildings are boxes standing on the ground, defined in the \code{[buildings]} section, one per key:
\code{name = x0, x1, height} (2D) or \code{name = x0, x1, y0, y1, height} (3D), in meters. Solid nodes have
zero velocity, no fuel and ambient values. Solid nodes with a fluid face neighbor (wall nodes) take
the average of their fluid neighbors for $T$ and the species: adiabatic, impermeable walls. With
the face-based projection the faces between solid and fluid carry no flux. Solids are opaque to
radiation at the temperature of their wall nodes, and the turbulence damping uses the distance to
the nearest building as well as the height above the ground.

Verification: a building in a 5~m\,s$^{-1}$ wind with all physics on and no fire keeps $T$ and
$\YO$ exactly at their ambient values with no flow inside the solid. Flow around a 10~m cube on a
1~m grid (10 cells per height, smooth power-law inflow) with the face-based projection: stable for
60~s, divergence error $1.6\times10^{-10}$~s$^{-1}$, mean upstream separation 0.8--1.5 heights ahead of
the front face and reattachment 3.5--3.7 heights behind the rear face. Wind-tunnel data for cubes
give about 1 and 1.6 heights (e.g.\ Martinuzzi \& Tropea 1993, channel flow); the longer wake is
attributed to the coarse resolution, the absence of inflow turbulence and the short averaging time.

\subsection{Terrain (immersed boundary)}
Nodes with height between the terrain height minus $\Delta z$ and the terrain height are \emph{cut
nodes}: velocity from the ground profile, zero vertical gradient of $T$ and species (from the two
nodes above). Nodes further below are \emph{dead nodes}: zero velocity and fuel, ambient species,
and the temperature of the surface node above them (adiabatic ground). The pressure equation is
solved over the whole domain.

%======================================================================
\section{Plan-view model}\label{sec:planview}
%======================================================================
\code{[model] type = planview} selects a reduced 2D model in the horizontal plane, of the
reaction--diffusion--convection type (Asensio \& Ferragut 2002). It has no vertical direction: the
plume, the indraft and the flame tilt are not resolved but parameterized, and the wind is
prescribed. It is a fire-spread model for the fuel-bed layer, not an atmosphere--fire model, and
complements the coupled model (for example for fast parameter sweeps or plan-view fire shapes).

State: $T(x,y)$, temperature of the fuel-bed layer (gas and fuel together), and $Y(x,y)$. With the
heat capacity of air plus fuel $C=\rho c_p+\rho_bc_sY$,
\begin{align}
  C\,\frac{\partial T}{\partial t} &= -\rho c_p\,\mathbf{U}\cdot\nabla T + \nabla\cdot\big(k(T)\nabla T\big)
    + \chi\,\rho_b H_C\Big(-\frac{\partial Y}{\partial t}\Big) - \frac{h_{top}}{Y_h}(T-T_\infty),\\
  \frac{\partial Y}{\partial t} &= -Y_f\,R(T)\,Y,
\end{align}
with $R$ of \eqref{eq:arrhenius} (integrated exactly over each step) and:
\begin{itemize}
  \item $\mathbf{U}$: mid-flame wind, speed $u_r(Y_h/z_r)^{\alpha_u}$ by default (the power law at
        the bed height; \code{[wind] factor} overrides the factor), direction \code{[wind] direction};
  \item $k(T)=\kappa+4\sigma\delta_{bed}T^3$: radiation as diffusion, justified here because the fuel
        bed is optically thick; $\delta_{bed}=4/(\alpha_s\sigma_s)$ (0.5~m for $\alpha_s=0.002$,
        $\sigma_s=4000$~m$^{-1}$) by default;
  \item $\chi$ (\code{[planview] heat\_retained}): fraction of the heat of combustion that stays in the bed;
  \item $h_{top}$ (\code{[planview] h\_top}): heat loss coefficient to the air above the bed.
\end{itemize}
Periodic boundaries in $x$ and $y$; Koren advection; the time integrators of the coupled model.
Ignition and fuel extent use the 3D keys (\code{[temperature] shape}, \code{x\_start}\dots\code{y\_end};
\code{[fuel] x\_start}\dots\code{y\_end}).

Status: $\chi$, $h_{top}$ and the wind factor carry the missing vertical physics and are not
calibrated. In a first test (F19 plot, line ignition, $\Delta x=0.5$~m) the behavior was
all-or-nothing: with $\chi=0.1$ the fire died, with $\chi=0.15$--$0.25$ it spread at about 3~m\,s$^{-1}$
but the bed temperature reached the 2500~K bound. Calibration of $\chi$, $h_{top}$, $T_{pc}$ and the
Arrhenius rate together, for example against the coupled model, is needed.

%======================================================================
\section{Numerical methods}\label{sec:numerics}
%======================================================================

\subsection{Grid and storage}
Uniform, collocated grid with $N_x$ ($N_y$) nodes in the periodic directions, of which the last
(equal to the first) is not stored, and $N_z$ nodes in $z$ including both walls. Spacings
$\Delta x=(x_{\max}-x_{\min})/(N_x-1)$, etc. All variables live at the nodes. Computations run in
PyTorch on CPU or GPU, by default in double precision.

\subsection{Spatial derivatives}
Periodic directions use central differences
$\partial_x\phi\approx(\phi_{i+1}-\phi_{i-1})/(2\Delta x)$ and
$\partial_{xx}\phi\approx(\phi_{i+1}-2\phi_i+\phi_{i-1})/\Delta x^2$.
In $z$ the same stencils are used in the interior, with second-order one-sided stencils at the
walls:
$\partial_z\phi_0\approx(-3\phi_0+4\phi_1-\phi_2)/(2\Delta z)$,
$\partial_{zz}\phi_0\approx(2\phi_0-5\phi_1+4\phi_2-\phi_3)/\Delta z^2$ (mirrored at the top).
The divergence in the projection uses the averaged half-node form
$\big(\tfrac12(\phi_{i+1}+\phi_i)-\tfrac12(\phi_i+\phi_{i-1})\big)/\Delta x$, which equals the
central difference.

\subsection{Advection}
\paragraph{Momentum: second-order upwind.}
\begin{equation}
  (a\,\partial_x\phi)_i \approx a_i^+\frac{3\phi_i-4\phi_{i-1}+\phi_{i-2}}{2\Delta x}
  + a_i^-\frac{-\phi_{i+2}+4\phi_{i+1}-3\phi_i}{2\Delta x},
\end{equation}
$a^\pm=\max/\min(a,0)$; in the two rows next to a wall the stencil that stays inside the domain is
used.

\paragraph{Temperature and species: Koren-limited third order (default).}
Face values are reconstructed from the upwind side,
\begin{equation}
  \phi_{i+1/2} = \phi_i + \tfrac12\,\psi(r_i)\,(\phi_{i+1}-\phi_i),\qquad
  r_i=\frac{\phi_i-\phi_{i-1}}{\phi_{i+1}-\phi_i},\qquad
  \psi(r)=\max\!\big(0,\min(2r,\tfrac{1+2r}{3},2)\big)
\end{equation}
for a non-negative face velocity $\tfrac12(a_i+a_{i+1})$ (mirrored otherwise), and
$(a\,\partial_x\phi)_i\approx a_i(\phi_{i+1/2}-\phi_{i-1/2})/\Delta x$. The scheme is third order
where $\phi$ is smooth and falls back to first-order upwind at extrema, so it creates no new
extrema for a uniform velocity. At the walls two ghost rows copy the boundary value. Alternatives:
\code{T\_advection = upwind} (second-order upwind) or \code{central} (the original scheme, which
oscillates at the fire front and with forward Euler is only stable because of the temperature bound).

\subsection{Time integration}
Explicit. The time step is adaptive by default (\code{[numerics] cfl}, target 0.4): each step takes
\begin{equation}
  \Delta t = \min\Big(\Delta t_{\max},\ \frac{\mathrm{CFL}}{\sum_j\max|u_j|/h_j},\
  \frac{0.25}{\nu_{\mathrm{eff}}\sum_j h_j^{-2}}\Big),\qquad \Delta t_{\max}=\frac{t_{\max}-t_{\min}}{N_t},
\end{equation}
with $\nu_{\mathrm{eff}}$ the largest molecular plus subgrid kinematic viscosity, and steps are
shortened to land exactly on the output times, which stay every $N_T\Delta t_{\max}$. In calm periods
$\Delta t=\Delta t_{\max}$ and the results equal those of a fixed step. With \code{cfl = 0} (and with
\code{--legacy}) the step is fixed. A fixed 1~ms step failed in an F19 calibration run when gusts
in the plume raised the CFL number from 0.43 to 0.73 within four seconds. The right-hand side $F(y)$ of all
equations except the pressure term is integrated with forward Euler (default), Heun's RK2 or
classical RK4; the projection is applied once per step, after the last stage, so the scheme is
first order in time overall. The radiation field $G$ is recomputed every \code{update\_every}
steps from the state at the start of the step and held fixed in between.

\subsection{Projection and pressure}
\paragraph{Face-based projection (default, \code{[pressure] projection = face}).}
Velocities are stored at the nodes, but mass conservation is imposed on face fluxes. Face velocities
are averages of the provisional node velocities, $F_{j,i+1/2}=\tfrac12(u_{j,i}+u_{j,i+1})$ (zero on
faces next to solids). The divergence of a node's control volume is the net flux through its faces,
$D_i=\sum_j(F_{j,i+1/2}-F_{j,i-1/2})/h_j$, where the ground node is a half cell with zero flux through the
ground. The projection
\begin{equation}
  F^{n+1}_{j,i+1/2} = F^*_{j,i+1/2} - \frac{\Delta t}{\rho_{i+1/2}}\,\frac{p_{i+1}-p_i}{h_j},\qquad
  \sum_j\frac{1}{h_j^2}\Big[\frac{p_{i+1}-p_i}{\rho_{i+1/2}}-\frac{p_i-p_{i-1}}{\rho_{i-1/2}}\Big]
  = \frac{D_i(F^*)-D^*_i}{\Delta t},
\end{equation}
with arithmetic face densities, makes the face fluxes satisfy the discrete constraint to the solver
tolerance (verified to about $10^{-11}$ for density ratios of 8 and with buildings). Faces next to
solid nodes are closed: zero flux through walls and a zero-gradient pressure condition, exactly.
The equation is solved by preconditioned conjugate gradients (symmetric with respect to the cell
volumes), preconditioned by the direct inverse of the constant-coefficient operator:
finite-difference eigenvalues $-(2\sin(\pi r/N)/h)^2$ in the periodic directions (FFT) and an
eigendecomposition of the vertical operator. With constant density and no solids one iteration
suffices; with buildings or strong density variations 10--30. The node velocities receive the
average of the corrections of their two faces (Rhie--Chow type), and the projected face fluxes
transport momentum, temperature and species in conservative form,
$\vu\cdot\nabla\phi\approx\sum_j\big(F_{j,i+1/2}\phi_{i+1/2}-F_{j,i-1/2}\phi_{i-1/2}\big)/h_j-\phi_iD_i$,
with Koren-limited face values $\phi_{i\pm1/2}$ upwinded with $F$.

\paragraph{Nodal projection (\code{projection = nodal}, original).}
Given the provisional state $y^*$ (velocity $\vu^*$), the pressure solves the variable-density
Poisson equation by fixed-point iteration:
\begin{equation}
  \nabla^2p^{(m+1)} = \frac{\rho^n}{\Delta t}\big(\nabla\cdot\vu^* - D^*\big)
  + \frac{\nabla\rho^n\cdot\nabla p^{(m)}}{\rho^n},
\end{equation}
with $D^*$ the divergence \eqref{eq:div} evaluated at $y^*$, until
$\max|p^{(m+1)}-p^{(m)}|<$\,\code{tol} or \code{max\_iter} iterations. Then
\begin{equation}
  \vu^{n+1} = \vu^* - \frac{\Delta t}{\rho^n}\nabla p.
\end{equation}

The nodal projection is not discretely consistent: the Laplacian of the Poisson solver is compact
(spectral in $x$, three points in $z$) while the divergence and gradient are central differences over
two cells, so the projected node velocities are not divergence-free for the central operators.
Errors of about 0.5~s$^{-1}$ were measured even away from obstacles in a flow around a cube, and the
flow around buildings went unstable; the face-based projection removes both problems.

\paragraph{Poisson solver (nodal projection).} FFT in the periodic directions (real FFT in $x$, complex in $y$). For
each horizontal wavenumber $k$ the vertical problem is tridiagonal: second-order differences in
$z$ with $\gamma_k=-2-(k\Delta z)^2$ on the diagonal, $p=0$ at the top row and a second-order
Neumann condition at the ground. Two direct methods solve all wavenumbers at once:
\begin{itemize}
  \item \code{thomas}: LU factors computed once; one forward and one backward sweep in $z$, each
        acting on a whole horizontal plane;
  \item \code{eig}: after eliminating the ground row the matrix is $M+\gamma_kI$ with the same $M$ for
        all $k$; one eigendecomposition of $M$ turns each solve into two dense products along $z$.
\end{itemize}
\code{auto} selects \code{eig} on GPUs (fewer, larger operations) and for small or medium grids on
CPU. Both agree to round-off.

\subsection{Radiation solver}
Gray discrete ordinates. Quadrature: \code{n\_polar} equal-solid-angle bands in
$\mu=\Omega_z$ (an even number, so no direction is horizontal) times \code{n\_azimuth} azimuths; in
2D the field does not depend on $y$, directions mirrored in $y$ give the same intensity and only
half the azimuths are kept with double weight, so a 2D run and a $y$-invariant 3D run give the same
$G$. Every direction has $\mu\ne0$, so the equation is integrated row by row in $z$ (upward
directions from the ground, downward ones from the top). Between rows the ray travels
$\Delta s=\Delta z/|\mu|$ and moves horizontally by $\Delta s\,\Omega_x$ (and $\Delta s\,\Omega_y$); the
upstream intensity is interpolated linearly (bilinearly in 3D) in the periodic directions, and
\begin{equation}
  I_k = I_{\text{upstream}}\,e^{-\kappa\Delta s} + \big(1-e^{-\kappa\Delta s}\big)\,I_b,
\end{equation}
exact for a uniform cell. All directions advance together, about 15 tensor operations per row.
Verification: exact equilibrium at uniform temperature; for an isothermal layer of optical
thickness $\tau$ the result matches $G=2\sigma T^4(1-E_2(\tau))$ within 0.3\% for $\tau=1$ with 16
polar bands (10\% low for $\tau=0.1$ with 6 bands).

\subsection{Stiff source terms}
\begin{itemize}
  \item Published reaction, \code{reaction = exact}:
        $\tilde R = (1-e^{-Y_fR\Delta t})/(Y_f\Delta t)$, so the explicit update gives the exact decay
        $Y e^{-Y_fR\Delta t}$ and $Y$ stays in $[0,1]$ (the explicit rate overshoots, and with RK2/RK4
        drives $Y$ negative within the stages).
  \item Pyrolysis is limited to the fuel present, $\dot m'''\le\rho_bY/\Delta t$.
  \item $\tau_{mix}\gg\Delta t$, so combustion is not stiff.
\end{itemize}

\subsection{Bounds and energy budget}
With \code{bound = True}: $T_\infty\le T\le T_{\max}$, $0\le Y\le1$, $T_\infty\le T_s\le T_{\max}$,
$0\le\Yf,\YO\le1$, $W\ge0$ after every step. The energy the temperature bounds add
($\rho c_p\Delta T$ summed over the clipped nodes) or remove is accumulated and reported against
the heat of combustion (\code{E\_clip\_low}, \code{E\_clip\_high}, \code{E\_combustion}). With the
original central scheme the lower bound added several times the combustion heat near the flame;
with the Koren scheme it is normally negligible.

\subsection{Log diagnostics}
Besides the extreme values, the log reports the total heat release rate and the gas temperature
weighted by the local heat release rate, $T_{flame}=\sum\dot q_c T/\sum\dot q_c$, which is comparable
with measured flame temperatures, unlike the single-cell maximum.

\subsection{Output}
2D runs save the full fields every $N_T$ steps; 3D runs save three planes (\code{xz\_*}, \code{xy\_*},
\code{yz\_*} through the head fire), the fuel left per column and the final full fields. Full
snapshots are written every \code{full\_interval} seconds (2D or 3D). Derived fields: heat release
rate (\code{HRR}), pyrolysis rate, convective heating of the fuel, $G$, net radiative heating of gas
and fuel, turbulent viscosity. Each run also stores its input file, all resolved parameters, git
information and a copy of the code.

%======================================================================
\section{Input parameters}\label{sec:params}
%======================================================================
Keys are given as \code{[section] key}. Values in a 2D file use \code{y} for the vertical
direction; a file is 3D when it defines \code{z\_min}, \code{z\_max} and \code{Nz}.

\begin{longtable}{p{4.6cm}p{2.6cm}p{7.6cm}}
\toprule
Key & Default & Meaning \\
\midrule
\endhead
\multicolumn{3}{l}{\emph{Domain and numerics}}\\
\code{[domain]} \code{x\_min}\dots\code{t\_max} & -- & domain and time interval \\
\code{[numerics] Nx, Ny, Nz} & -- & nodes per direction \\
\code{[numerics] Nt, NT} & -- & time steps; output every \code{NT} steps \\
\code{[numerics] time\_method} & \code{euler} & \code{euler}, \code{RK2}, \code{RK4} \\
\code{[numerics] cfl} & 0.4 & adaptive step target; 0: fixed $\Delta t=(t_{\max}-t_{\min})/N_t$ \\
\code{[numerics] T\_advection} & \code{koren} & \code{koren}, \code{upwind}, \code{central} \\
\code{[numerics] bound} & \code{True} & apply the bounds of the solution \\
\code{[pressure] tol, max\_iter} & $10^{-10}$, 50 & fixed point of the pressure \\
\code{[pressure] poisson} & \code{auto} & \code{auto}, \code{thomas}, \code{eig} (nodal projection) \\
\code{[pressure] projection} & \code{face} & \code{face} or \code{nodal} \\
\code{[buildings]} \code{name} & -- & building box (see Section~\ref{sec:ibc}) \\
\code{[density] constant} & \code{False} & constant density (incompressible) \\
\multicolumn{3}{l}{\emph{Gas and turbulence}}\\
$c_p$, $\rho_\infty$, $\mu$, $\kappa$ & 1007, 1.225, $1.802\times10^{-5}$, 0.02476 & SI units (not in the input file) \\
\code{[gas] cp\_model} & \code{constant} & \code{constant}, \code{variable} ($c_p(T)$ of air), \code{mixture} \\
\code{[gas] transport} & \code{constant} & \code{constant} or \code{sutherland} \\
\code{[temperature] T\_inf} & 288.15 K & ambient temperature \\
\code{[turbulence] enabled} & \code{True} & Smagorinsky model \\
\code{[turbulence] C\_s} & 0.2 & Smagorinsky constant \\
\code{[turbulence] filter\_width} & grid & fixed filter width (m) \\
$Pr$ & 0.7323 & Prandtl number (also turbulent) \\
\multicolumn{3}{l}{\emph{Fuel}}\\
\code{[fuel] Y\_h} & 0.51 m & fuel bed height \\
\code{[fuel] C\_d, alpha\_s, sigma\_s} & 0.15, 0.002, 4000 m$^{-1}$ & drag coefficient, solid fraction, surface-to-volume \\
\code{[fuel] x\_start}\dots\code{y\_end} & whole domain & horizontal extent of the bed \\
\code{[fuel] H\_C} & $21\times10^6$ J\,kg$^{-1}$ & heat of combustion \\
\code{[fuel] A, T\_act, Y\_f} & $10^9$ s$^{-1}$, 18040.85 K, 70 & Arrhenius rate (published model, \code{arrhenius} pyrolysis) \\
\code{[temperature] T\_pc} & 523 K & phase-change temperature \\
\code{[fuel] reaction} & \code{exact} & \code{exact} or \code{explicit} \\
\code{[temperature] sink, h\_c} & \code{surface}, 1.42 & convective exchange (published model), coefficient \\
\code{[fuel] solid\_temperature} & \code{True} & extended model \\
\code{[fuel] rho\_bulk} & 0.6144 kg\,m$^{-3}$ & dry bulk density (FDS Case F19) \\
\code{[fuel] c\_s} & 1500 J\,kg$^{-1}$K$^{-1}$ & specific heat of dry fuel \\
\code{[fuel] H\_pyr} & $4.16\times10^5$ J\,kg$^{-1}$ & heat of pyrolysis \\
\code{[fuel] moisture} & 0 & moisture content (F19: 0.058) \\
\code{[fuel] pyrolysis} & \code{linear} & \code{linear} or \code{arrhenius} \\
\code{[fuel] T\_pyr, T\_pyr\_end} & 400 K, 500 K & pyrolysis range \\
$c_w$, $L_v$, $T_{boil}$ & 4184, $2.259\times10^6$, 373.15 & water properties (SI) \\
\multicolumn{3}{l}{\emph{Combustion and radiation}}\\
\code{[combustion] s\_O2} & 1.185 & oxygen per kg of fuel gas \\
\code{[combustion] tau\_mix} & $\sqrt{2\Delta/|g|}$ & mixing time (s) \\
\code{[combustion] Y\_O2\_amb} & 0.232 & ambient oxygen \\
\code{[radiation] model} & \code{transport} & \code{transport} or \code{diffusion} \\
\code{[radiation] kappa\_gas} & 0.1 m$^{-1}$ & gray gas absorption coefficient \\
\code{[radiation] n\_polar, n\_azimuth} & 12, 16 & angular quadrature (even) \\
\code{[radiation] update\_every} & 5 & steps between radiation solves \\
\code{[temperature] delta} & 0.1 m & optical path length (diffusion radiation) \\
\multicolumn{3}{l}{\emph{Initial conditions}}\\
\code{[wind] u\_r, z\_r, shape} & 4.8 m\,s$^{-1}$, 2 m, power law & reference wind and profile \\
\code{[wind] z\_0, L} & 0.05 m, neutral & roughness and Obukhov length (\code{monin-obukhov}) \\
\code{[wind] forcing, forcing\_time} & \code{none}, 1 s & wind forcing (\code{profile}) and its relaxation time \\
\code{[wall] model, z\_0} & \code{none}, wind $z_0$ & ground wall model (\code{log}) and its roughness \\
\code{[temperature] T\_hot, shape} & 788.15 K, \code{gaussian} & ignition temperature and shape \\
\code{[temperature] x\_start}\dots\code{z\_end} & 0, 6, -, -, 0, 1 m & ignition extent \\
\code{[temperature] t\_source} & $-1$ (off) & sustained source until this time \\
\code{[temperature] T\_min, T\_max} & $T_\infty$, 2500 K & temperature bounds \\
\code{[topography] shape} & \code{flat} & \code{flat}, \code{hill}, \code{slope} (2D) \\
\multicolumn{3}{l}{\emph{Output}}\\
\code{[output] z\_slice, y\_slice} & 1 m, center & planes saved in 3D \\
\code{[output] full\_interval} & 0 (off) & seconds between full snapshots \\
\code{[output] full\_dtype} & \code{float32} & precision of the snapshots \\
\bottomrule
\end{longtable}

%======================================================================
\section{Known limitations and open issues}
%======================================================================
\begin{itemize}
  \item Not calibrated: $\kappa_g$, $H_C$, $\tau_{mix}$ and $\rho_b$ control the heat release; peak
        gas temperatures of about 2200--2400~K suggest too high a heat release (constant $c_p$,
        $H_C$ of the solid rather than of the pyrolysis gas, weak flame radiation).
  \item Constant $c_p$; no dissociation; one gray absorption coefficient for the gas, independent of
        soot or products.
  \item Water vapor from drying and combustion products are not transported.
  \item The ignition is a hot-gas pulse at $t=0$, not the timed burner of the FDS F19 setup.
  \item Periodic horizontal boundaries instead of open ones; without \code{[wind] forcing} the
        near-ground wind decays during a run, and without the wall model the no-slip ground
        overestimates the friction on coarse grids (Section~\ref{sec:ibc}). Building walls have no
        wall model.
  \item The intermittent activity of the lower temperature bound seen on the 1.17~m F19 grid was
        traced to the second-order zero-gradient condition at the ground (Section~\ref{sec:ibc}) and
        fixed: over the first 5~s of that case the energy added by the bound fell from 16\% to 1.6\% of
        the heat of combustion, and the worst undershoot from about 250~K to under 2~K.
\end{itemize}

\section*{References}
\begin{itemize}
  \item San Mart\'in, D., \& Torres, C.~E. (2025). Simplified coupled atmosphere--fire model for
        simulation of 2D wildland fires. \emph{Computer Physics Communications}.
        \url{https://doi.org/10.1016/j.cpc.2025.109746}
  \item Mell, W., Jenkins, M.~A., Gould, J., \& Cheney, P. (2007). A physics-based approach to
        modeling grassland fires. \emph{International Journal of Wildland Fire}, 16, 1--22.
  \item McGrattan, K., et al. \emph{Fire Dynamics Simulator, Technical Reference Guide}. NIST.
  \item Wrede, H., et al. (2025). Fire-X: Extinguishing fire with stoichiometric heat release.
        \emph{ACM Transactions on Graphics}, 44(6), 269.
  \item Martinuzzi, R., \& Tropea, C. (1993). The flow around surface-mounted, prismatic obstacles
        placed in a fully developed channel flow. \emph{Journal of Fluids Engineering}, 115, 85--92.
  \item Rhie, C.~M., \& Chow, W.~L. (1983). Numerical study of the turbulent flow past an airfoil
        with trailing edge separation. \emph{AIAA Journal}, 21, 1525--1532.
  \item Paulson, C.~A. (1970). The mathematical representation of wind speed and temperature
        profiles in the unstable atmospheric surface layer. \emph{Journal of Applied Meteorology}, 9,
        857--861.
  \item Koren, B. (1993). A robust upwind discretization method for advection, diffusion and
        source terms. In \emph{Numerical Methods for Advection--Diffusion Problems}, Vieweg.
\end{itemize}

\end{document}
